% w12_s02_guia_docente_facilitacion_diferenciada.tex
% IMT-221 — Semana 12, Sesión 2 (21/10/2026)
% Artefacto 4: Guía docente breve de facilitación diferenciada (105 min)
\PassOptionsToPackage{table,xcdraw}{xcolor}
\documentclass[10pt, a4paper]{article}

\usepackage[T1]{fontenc}
\usepackage[utf8]{inputenc}
\usepackage{tgpagella}
\usepackage{tgcursor}
\usepackage[spanish, es-noshorthands]{babel}
\usepackage{amsmath, amssymb}
\usepackage{booktabs}
\usepackage{tabularx}
\usepackage{xcolor}
\usepackage{tikz}
\usetikzlibrary{shapes.geometric, arrows.meta, positioning}
\usepackage[american]{circuitikz}
\usepackage[most]{tcolorbox}
\usepackage[margin=1.6cm]{geometry}
\usepackage{fancyhdr}

\setlength{\headheight}{16pt}
\setlength{\emergencystretch}{2em}
\definecolor{ucbblue}{RGB}{0, 51, 102}
\definecolor{neutralbg}{RGB}{247, 248, 250}
\definecolor{warnbg}{RGB}{254, 246, 246}

\newcolumntype{P}[1]{>{\raggedright\arraybackslash}p{#1}}
\newcolumntype{Y}{>{\raggedright\arraybackslash}X}

\pagestyle{fancy}
\fancyhf{}
\lhead{\textbf{IMT-221}}
\chead{Guía Docente de Facilitación Diferenciada (105 min)}
\rhead{Semana 12 - Sesión 2}
\lfoot{IMT-221 — Circuitos Electrónicos II}
\rfoot{Página \thepage}
\renewcommand{\headrulewidth}{0.4pt}

\begin{document}

\begin{center}
    {\Large\textbf{\textcolor{ucbblue}{Guía Docente de Facilitación Diferenciada y Auditoría Técnica (105 min)}}} \\[0.1cm]
    {\footnotesize\textbf{Semana 12 — Sesión 2 (21/10/2026) · Estado del Insumo: Planificación Propuesta para Ejecución Controlada}}
\end{center}

\section*{1. Guion Operativo de la Sesión (105 Minutos — Prioridad: Seguridad y Dictámenes Reales)}
\begin{table}[htbp]
    \centering
    \footnotesize
    \renewcommand{\arraystretch}{1.24}
    \begin{tabularx}{\textwidth}{|c|c|Y|P{4.8cm}|}
    \hline
    \rowcolor[HTML]{EFEFEF} 
    \textbf{Bloque} & \textbf{Min} & \textbf{Acción Docente en Aula / Laboratorio} & \textbf{Artefacto / Evidencia Generada} \\ \hline
    1. Contrato e informes & 8 & Verificar estado real de recepción Lab 3 (plazo 18/10 vencido), recordar plazo Lab 4 (25/10) y fijar regla de trazabilidad. & Registro administrativo inicial. \\ \hline
    2. Ficha mínima & 15 & Cada grupo completa su Ficha Única de Evidencia declarando archivos exactos o \texttt{NO MEDIDO}. & Artefacto 1 completado por grupo. \\ \hline
    3. Revisión Gate A & 25 & Auditoría rápida de cálculos, LTSpice, datasheet y fuente; emitir dictamen (\texttt{READY}/\texttt{MINOR}/\texttt{NOT READY}) y un bloqueo causal. & Acta de Gate A firmada (Artefacto 1). \\ \hline
    4. Trabajo diferenciado & 25 & \texttt{READY}: montaje aislado tras pre-power. \texttt{MINOR}: ajuste puntual y recheck. \texttt{NOT READY}: reparación de diseño/SPICE sin energizar. & Medición incremental o corrección de diseño. \\ \hline
    5. Factibilidad Hito 2 & 12 & Llenar matriz colectiva de avance y evaluar viabilidad técnica de la defensa propuesta del 28/10 (\texttt{GO/NO-GO}). & Artefacto 2 (Matriz consolidada). \\ \hline
    6. Mapa curricular & 15 & Presentar ruta de conexión Hito 2 $\rightarrow$ Unidad 3 (filtros activos) $\rightarrow$ Hito 3 sin dar por aprobados cambios de fechas. & Artefacto 3 (Resumen estudiantil). \\ \hline
    7. Cierre individual & 5 & Muestreo oral individual (1 pregunta discriminante por estudiante) y asignación de tarea exacta por grupo. & Acta de cierre y spot-check. \\ \hline
    \end{tabularx}
\end{table}
\noindent\textit{Regla de gestión de tiempo:} Si falta tiempo en la sesión, reduzca las discusiones plenarias, \textbf{nunca los controles de seguridad eléctrica pre-energizado}.

\section*{2. Verificación Matemática Rápida del Diseño del Grupo (Bloque 3 y Ejemplo 1)}
Al inspeccionar el esquema existente del grupo, verifique los siguientes puntos críticos y relaciones algebraicas paso a paso:

\begin{center}
    \begin{circuitikz}[scale=0.74, transform shape, thick]
        \draw (0,2.2) node[op amp, noinv input up] (A1) {$A_1$};
        \draw (0,-2.2) node[op amp, noinv input down] (A2) {$A_2$};
        \draw (A1.+) to[short, -o] ++(-1.2,0) node[left] {$V_1 = V_{cm} - \frac{V_d}{2}$};
        \draw (A2.+) to[short, -o] ++(-1.2,0) node[left] {$V_2 = V_{cm} + \frac{V_d}{2}$};
        \draw (A1.-) -- ++(0,-0.9) coordinate(n1);
        \draw (A2.-) -- ++(0,0.9) coordinate(n2);
        \draw (n1) to[R, l={$R_G$}, i>^=$I_G$] (n2);
        \draw (n1) to[R, l={$R$}] (n1 -| A1.out) -- (A1.out);
        \draw (n2) to[R, l_={$R$}] (n2 -| A2.out) -- (A2.out);
        \draw (A1.out) to[short, *-] ++(1.2,0) coordinate(vo1) node[above] {$V_{o1}$ (Revisar swing)};
        \draw (A2.out) to[short, *-] ++(1.2,0) coordinate(vo2) node[below] {$V_{o2}$ (Revisar swing)};
        \draw (6.6,0) node[op amp] (A3) {$A_3$};
        \draw (vo1) -- (vo1 |- A3.-) to[R, l={$R_1$}] (A3.-);
        \draw (vo2) -- (vo2 |- A3.+) to[R, l={$R_3$}] (A3.+);
        \draw (A3.-) -- ++(0,1.4) coordinate(f3) to[R, l={$R_2$}] (f3 -| A3.out) -- (A3.out);
        \draw (A3.+) to[R, l={$R_4$}] ++(0,-1.4) node[ground]{} node[right=0.15cm] {$V_{REF}$};
        \draw (A3.out) to[short, *-o] ++(1.1,0) node[right] {$V_o = k(V_{o2}-V_{o1})$};
    \end{circuitikz}
\end{center}

\textbf{Derivación explícita que el estudiante debe poder defender:}
\begin{align}
    V_d &= V_2 - V_1, \qquad V_{cm} = \frac{V_1 + V_2}{2} \quad \Longleftrightarrow \quad V_1 = V_{cm} - \frac{V_d}{2}, \qquad V_2 = V_{cm} + \frac{V_d}{2} \\
    I_G &= \frac{V_1 - V_2}{R_G} \implies V_{o1} = V_1 + \frac{R}{R_G}(V_1 - V_2), \qquad V_{o2} = V_2 + \frac{R}{R_G}(V_2 - V_1) \\
    V_{o2} - V_{o1} &= (V_2 - V_1) + \frac{2R}{R_G}(V_2 - V_1) = \left(1 + \frac{2R}{R_G}\right)V_d \\
    \text{Con } k &= \frac{R_2}{R_1} = \frac{R_4}{R_3}: \quad V_o = k(V_{o2} - V_{o1}) \implies \boxed{A_d = k\left(1 + \frac{2R}{R_G}\right)} \implies \boxed{R_G = \frac{2R}{(A_d^\star / k) - 1}} \quad (A_d^\star > k) \\
    \text{Con } R_G^\prime &\text{ real disponible:} \quad \boxed{A_{d,\mathrm{nom}} = k\left(1 + \frac{2R}{R_G^\prime}\right)}, \qquad \boxed{V_{o,\mathrm{pred}} = A_{d,\mathrm{nom}}V_d} \\
    V_o &= A_d V_d + A_c V_{cm}, \qquad \boxed{CMRR = \left|\frac{A_d}{A_c}\right|}, \qquad \boxed{CMRR_{\mathrm{dB}} = 20\log_{10}\left|\frac{A_d}{A_c}\right|}
\end{align}

\textbf{Ejemplo Trabajado 1 — Auditoría Completa de una Ficha Gate A (Datos Hipotéticos):} \\
Un grupo presenta medición de Hito 1 con $V_{d,\max} = 40\,\text{mV} = 0.040\,\text{V}$, fija $V_{o,\mathrm{target}} = 0.80\,\text{V}$, elige $k=1$, $R=10\,\text{k}\Omega$ y declara usar $R_G^\prime = 1.0\,\text{k}\Omega$, pero no adjunta prueba común en SPICE ni datasheet del integrado:
\begin{align*}
    1)\; A_d^\star &= \frac{0.80\,\text{V}}{0.040\,\text{V}} = 20, \qquad 2)\; R_G = \frac{2(10\,\text{k}\Omega)}{20 - 1} = \frac{20\,\text{k}\Omega}{19} \approx 1.0526\,\text{k}\Omega \approx 1.053\,\text{k}\Omega \\
    3)\; A_{d,\mathrm{nom}} &= 1 + \frac{2(10\,\text{k}\Omega)}{1.0\,\text{k}\Omega} = 21, \qquad 4)\; V_{o,\mathrm{pred},\max} = 21(0.040\,\text{V}) = 0.84\,\text{V} \quad (\text{condicional a no saturar})
\end{align*}
\textit{Dictamen docente del ejemplo:} \textbf{\texttt{NOT READY}} (No clasificar como brecha documental menor). Faltan datos críticos de simulación en modo común, identificación del IC físico y verificación de márgenes eléctricos en el datasheet.

\section*{3. Seguridad Eléctrica Pre-Power, Rieles Efectivos y Tierra de Instrumentos (Bloque 5)}
\begin{center}
    \begin{circuitikz}[scale=0.80, transform shape, thick, >=Stealth]
        % Divisor y Buffer en el lado izquierdo (x de 0 a 3.6)
        \draw (0,2.2) node[vcc]{$V_S = +18\,\text{V}$} to[R, l={$R_a$}] (0,0.9) coordinate(vm) to[R, l={$R_b$}] (0,-0.4) node[ground]{$0\,\text{V}_{\text{físico}}$};
        \draw (vm) to[short, *-] ++(0.8,0) node[op amp, noinv input up, anchor=+, scale=0.72] (buf) {};
        \draw (buf.-) -- ++(0,-0.75) -| (buf.out);
        \draw (buf.out) to[short, *-o] ++(1.1,0) coordinate(vref);
        
        % Etiqueta del nodo V_REF colocada arriba/abajo para no invadir la derecha
        \node[above=3pt, font=\footnotesize\bfseries, text=ucbblue] at (vref) {$V_{REF} = 8.7\,\text{V}$};
        \node[below=3pt, font=\scriptsize] at (vref) {($0\,\text{V}_{\text{circuito}}$)};

        % Flecha de advertencia y cuadro desplazado con anchor=west en x = 5.6 (cero solapamiento)
        \draw[<->, dashed, red!80!black, thick] (vref) ++(0.25,0) -- node[above, font=\scriptsize\bfseries, text=red!80!black] {¿Pinza GND?} (5.5, 0.9);
        \node[draw=red!80!black, fill=warnbg, rounded corners=2pt, text width=6.6cm, align=center, font=\scriptsize, inner sep=4pt, anchor=west] at (5.6, 0.9) {\textbf{¿Tierra del Osciloscopio / Generador conectada a Tierra de Protección?}\\$\mathbf{?}$ \textbf{NO CONECTAR pinza de masa a $V_{REF}$ sin verificar aislamiento en manual y con supervisión docente}};
    \end{circuitikz}
\end{center}

\textbf{Ejemplo Trabajado 3 — Cálculo de Rieles Efectivos en Referencia Virtual:} \\
Si la fuente física entrega de $0\,\text{V}$ a $+18.0\,\text{V}$ y la salida del buffer se mide en $V_{REF} = 8.7\,\text{V}$ respecto al terminal de $0\,\text{V}$:
\begin{equation}
    V_{+,eff} = V_S - V_{REF} = 18.0\,\text{V} - 8.7\,\text{V} = \boxed{+9.3\,\text{V}}, \qquad V_{-,eff} = 0.0\,\text{V} - 8.7\,\text{V} = \boxed{-8.7\,\text{V}}
\end{equation}
\textbf{Regla de Decisión Pre-Power:} LM741/UA741, LM358 y TL074 \textbf{no son intercambiables pin a pin}. Si falta cualquier dato de seguridad eléctrica (pinout exacto, rieles medidos sin IC, compatibilidad de masa del instrumento), el dictamen obligatorio es \textbf{\texttt{NO ENERGIZAR}}.

\section*{4. Preguntas Rápidas de Discriminación y Diagnóstico Individual (Ejercicios 1 a 8)}
\begin{table}[htbp]
    \centering
    \scriptsize
    \renewcommand{\arraystretch}{1.28}
    \begin{tabularx}{\textwidth}{|c|P{6.2cm}|Y|}
    \hline
    \rowcolor[HTML]{EFEFEF} 
    \textbf{\#} & \textbf{Pregunta / Ejercicio de Diagnóstico en Clase} & \textbf{Respuesta / Criterio Docente Esperado} \\ \hline
    1 & Muestra exactamente de dónde sale tu $V_{d,\max}$, su condición y cómo lo conviertes en $A_d^\star$. & Debe abrir el registro de Hito 1 y mostrar $A_d^\star = V_{o,\mathrm{target}}/V_{d,\max}$. Si no existe, definir la medición faltante (no inventar). \\ \hline
    2 & Para $A_d^\star = 11$, $R = 10\,\text{k}\Omega$ y $k = 1$, calcula $R_G$ y explica por qué no fija resistores para todos. & $R_G = \frac{2(10\,\text{k}\Omega)}{11 - 1} = \frac{20\,\text{k}\Omega}{10} = \boxed{2.0\,\text{k}\Omega}$. Cada grupo tiene su propio $V_{d,\max}$ de Hito 1 y su propio inventario real. \\ \hline
    3 & Con $V_1 = V_{cm} - V_d/2$ y $V_2 = V_{cm} + V_d/2$ en LTSpice, ¿cómo compruebas linealidad? & Evaluando al menos dos niveles de $V_d$, verificando que $\Delta V_o/\Delta V_d \approx A_{d,\mathrm{nom}}$ y que $V_{o1}, V_{o2}, V_o$ no se recorten contra los rieles. \\ \hline
    4 & ¿Qué cambia físicamente si reemplazas un LM741 por un TL074 en tu protoboard? & Cambia el encapsulado (DIP-8 a DIP-14), pinout completo, rango de modo común, swing y requisitos de fuente; exige redibujar y recablear. \\ \hline
    5 & Identifica el nodo homólogo antes y después de insertar el In-Amp para evaluar efecto de carga. & El nodo de salida diferencial del puente de Wheatstone bajo idéntica excitación, midiendo la tensión preservada antes de la ganancia activa. \\ \hline
    6 & «Tengo esquema y LTSpice, pero desconozco el IC exacto y la fuente física»: ¿cuál es el estado? & \textbf{\texttt{NOT READY}} para energizar. Faltan identificación del chip, datasheet, pinout, medición de rieles reales y verificación de tierras. \\ \hline
    7 & Justifica por qué medir $A_d = 50$ no demuestra un alto CMRR físico. & Porque $CMRR = |A_d/A_c|$; si $A_c$ es grande por desajuste de razones ($R_2/R_1 \neq R_4/R_3$), el rechazo común será bajo aunque $A_d$ sea 50. \\ \hline
    8 & Distingue una actividad técnica de Unidad 3 de la evidencia de Hito 3 y del examen individual. & Unidad 3 exige diseñar/analizar filtros activos y estabilidad; Hito 3 exige un producto integrado diferenciado; el examen evalúa dominio individual. \\ \hline
    \end{tabularx}
\end{table}

\section*{5. Mapa de Reutilización de Materiales Existentes y Acta de Cierre}
\begin{table}[htbp]
    \centering
    \scriptsize
    \renewcommand{\arraystretch}{1.22}
    \begin{tabularx}{\textwidth}{|P{6.8cm}|Y|}
    \hline
    \rowcolor[HTML]{EFEFEF}
    \textbf{Archivo Existente (Reutilizar sin regenerar)} & \textbf{Función Específica en la Sesión del 21/10/2026} \\ \hline
    {\ttfamily w12\_s01\_guia\_comprensiva\_gate\_A\_hito\_2.pdf} & Consulta conceptual de pasos y ejemplos de Gate A. \\ \hline
    {\ttfamily w12\_s01\_taller\_diseno\_depurado.pdf} & Reparación modular puntual para grupos en estado {\ttfamily NOT READY}. \\ \hline
    {\ttfamily w12\_s01\_anexo\_verificacion\_ltspice\_gate\_A.pdf} & Auditoría de pruebas diferencial, común y headroom en SPICE. \\ \hline
    {\ttfamily w12\_s01\_guia\_construccion\_fisica\_hito\_2.pdf} & Protocolo de montaje aislado para grupos {\ttfamily READY}. \\ \hline
    {\ttfamily w12\_s01\_registro\_evidencia\_fisica\_hito\_2.pdf} & Registro estandarizado de rieles y mediciones incrementales. \\ \hline
    {\ttfamily w12\_s01\_anexo\_docente\_troubleshooting.pdf} y {\ttfamily w12\_s01\_banco\_preguntas\_individuales.pdf} & Diagnóstico sistemático de fallas y muestreo oral individual. \\ \hline
    \end{tabularx}
\end{table}

\begin{tcolorbox}[colback=neutralbg, colframe=black, boxsep=2pt, top=3pt, bottom=3pt, title=\textbf{Síntesis Docente de Cierre de Sesión (21/10/2026)}]
    \footnotesize\raggedright
    \textbf{Grupos evaluados:} \rule{2.0cm}{0.4pt} \hfill \textbf{READY:} \rule{2.0cm}{0.4pt} \hfill \textbf{MINOR:} \rule{2.0cm}{0.4pt} \hfill \textbf{NOT READY:} \rule{2.0cm}{0.4pt} \\[0.15cm]
    \textbf{Etapas aisladas medidas hoy:} \rule{2.2cm}{0.4pt} \hfill \textbf{Recomendación 28/10:} $\square$ \texttt{GO} \quad $\square$ \texttt{NO-GO} \\[0.15cm]
    \textbf{Observaciones críticas de seguridad / instrumentación / inventario:} \hrulefill \\[0.15cm]
    \textbf{Firma del Docente:} \rule{5cm}{0.4pt} \hfill \textbf{Fecha y Hora de Cierre:} \rule{4cm}{0.4pt}
\end{tcolorbox}

\end{document}
